\providecommand{\dd}{\,\mathrm{d}}
\AtBeginDocument{\setlength{\abovedisplayskip}{4pt plus 1pt minus 1pt}
\setlength{\belowdisplayskip}{4pt plus 1pt minus 1pt}
\setlength{\abovedisplayshortskip}{3pt plus 1pt}
\setlength{\belowdisplayshortskip}{3pt plus 1pt}
\setlength{\jot}{2pt}
}
\tcbset{handout base/.append style={before skip=2.5pt,after skip=2.5pt,before upper={\RaggedRight}}}

\HandoutSetup{NE 630}{13}{The Multigroup Method}
  {FNRP Ch.~3 (multigroup equations); review Lesson 12}
\begin{document}
\MakeHandoutTitle
\textcolor{KSUDeepPurple}{\textbf{Primary objective.}}
Construct multigroup cross sections and solve neutron balances
for an infinite, homogeneous system.
\par\vspace{5pt}

\begin{handoutcolumns}
\HandoutSection{1}{Replace assumed flux ratios}
Lesson 12 supplied normalized shapes and \emph{assumed} the
flux in each range. Now determine those fluxes from coupled balances.
Number groups from \textbf{high to low energy}:
\[
 E_0>E_1>\cdots>E_G,
 \qquad \phi_g=\int_{E_g}^{E_{g-1}}\phi(E)\dd E.
\]
The complete domain has $E_0=\infty$ and $E_G=0$.
For our three-group examples, use the following finite bounds (eV):
\[
 (E_0,E_1,E_2,E_3)=(10^7,10^5,1,10^{-3}).
\]
Thus $(\phi_1,\phi_2,\phi_3)=(\phi_F,\phi_I,\phi_T)$.
Finite cutoffs require attention to transfers outside the retained range.

\HandoutSection{2}{Preserve each group reaction rate}
For $\phi_g>0$, define
\begin{fixedbox}{Group averages}
\begin{align*}
 \Sigma_{x,g}&=
 \frac{\int_{E_g}^{E_{g-1}}\Sigma_x(E)\phi(E)\dd E}{\phi_g},\\
 (\nu\Sigma_f)_g&=
 \frac{\int_{E_g}^{E_{g-1}}\nu(E)\Sigma_f(E)\phi(E)\dd E}{\phi_g}.
\end{align*}
\end{fixedbox}
Use a normalized within-group shape to prepare these averages;
then solve for $\phi_g$ (Lesson 12). The group sources are
\[
 \chi_g=\int_{E_g}^{E_{g-1}}\chi(E)\dd E,
 \qquad s_g=\int_{E_g}^{E_{g-1}}s_{\mathrm{ext}}'''(E)\dd E.
\]

\HandoutSection{3}{Keep the scattering direction straight}
\begin{fixedbox}{Destination $g$ $\gets$ incident group $g'$}
\[
 \Sigma_{s,g\gets g'}=
 \frac{\displaystyle
 \int_{E_{g'}}^{E_{g'-1}}\!\phi(E')
 \int_{E_g}^{E_{g-1}}\!\Sigma_s(E'\!\to E)\dd E\dd E'}
 {\phi_{g'}}.
\]
\end{fixedbox}
The denominator is the \textbf{incident-group flux}.
Multiplying by $\phi_{g'}$ recovers the rate of transfer into $g$.
\begin{checkpoint}
With this numbering, downscatter means
$g\ \MathBlank[0.35in]{}\ g'$. In a scattering matrix whose
rows are destinations, are the forced zero (upscatter) entries
above or below the diagonal?
\RuledLines{2}
\end{checkpoint}

\switchcolumn
\RaggedRight
\HandoutSection{4}{One balance per destination group}
Integrate the spectrum equation over $E_g\le E\le E_{g-1}$:
\begin{fixedbox}{Infinite-medium multigroup balance}
\begin{align*}
 \Sigma_{t,g}\phi_g
 &=\sum_{g'=1}^G\Sigma_{s,g\gets g'}\phi_{g'}\\
 &\quad+\chi_g\sum_{g'=1}^G(\nu\Sigma_f)_{g'}\phi_{g'}+s_g,
 \qquad g=1,\ldots,G.
\end{align*}
\end{fixedbox}
The fission term sums production over \emph{all incident groups};
$\chi_g$ then distributes the newborn neutrons to destination $g$.
Move self-scattering to the left by defining \textbf{removal}:
\[
 \boxed{\Sigma_{r,g}=\Sigma_{t,g}-\Sigma_{s,g\gets g}}.
\]
With number conservation and complete energy coverage,
\[
 \Sigma_{r,g}=\Sigma_{a,g}
      +\sum_{h\ne g}\Sigma_{s,h\gets g}.
\]
Removal includes absorption \emph{and} scattering out of a group;
it is not generally the same as absorption.

\HandoutSection{5}{Three-group water tank}
Approximate the ${}^{252}$Cf source as $(s_1,s_2,s_3)=(10,0,0)$
in $\mathrm{cm^{-3}\,s^{-1}}$. Ignore multiplication, leakage,
intergroup upscatter, and transfers outside the finite energy range.
\begin{boardbox}{Solve from fast to thermal}
Write each group balance and isolate its flux:
\begin{align*}
 \phi_1&=\MathBlank[1.30in]{},\\
 \phi_2&=\MathBlank[1.75in]{},\\
 \phi_3&=\frac{\Sigma_{s,3\gets1}\phi_1
                  +\Sigma_{s,3\gets2}\phi_2}{\Sigma_{r,3}}.
\end{align*}
\end{boardbox}
Unlike the Lesson 12 example, the flux ratios are now outputs.
\begin{checkpoint}[Check the complete balance]
Add the three equations. Under the stated assumptions, show that
\[
 \sum_{g=1}^3\Sigma_{a,g}\phi_g=10\unit{cm^{-3}\,s^{-1}}.
\]
Why should removal cross sections \emph{not} replace absorption
cross sections in this check?
\RuledLines{2}
\end{checkpoint}
\end{handoutcolumns}

\newpage
\begin{handoutcolumns}
\HandoutSection{6}{Use the scattering kinematics}
For elastic scattering isotropic in the CM system from a
stationary target of mass ratio $A$, let
$\alpha=[(A-1)/(A+1)]^2$. Then
\[
 p(E'\!\to E)=
 \begin{cases}
  1/[(1-\alpha)E'],&\alpha E'\le E\le E',\\
  0,&\text{otherwise}.
 \end{cases}
\]
This kernel cannot upscatter; it is not a thermal free-gas kernel.
The allowed outgoing interval and destination group overlap over
\begin{align*}
 \Delta E_g(E')=\Big[&\min(E_{g-1},E')\\
                  &-\max(E_g,\alpha E')\Big]_+,
 \qquad [z]_+=\max(z,0).
\end{align*}
Thus the probability of landing in group $g$ is
\[
 P_g(E')=\frac{\Delta E_g(E')}{(1-\alpha)E'}.
\]
\begin{fixedbox}{Incident-flux-weighted transfer}
\[
 \boxed{\sigma_{s,g\gets g'}=
 \frac{\displaystyle\int_{E_{g'}}^{E_{g'-1}}
      \sigma_s(E')P_g(E')\phi(E')\dd E'}
      {\displaystyle\int_{E_{g'}}^{E_{g'-1}}\phi(E')\dd E'}}.
\]
\end{fixedbox}
For a mixture, compute each nuclide's transfer with its own
$\alpha_i$, then form $\Sigma_{s,g\gets g'}=\sum_iN_i\sigma^i_{s,g\gets g'}$.
\begin{checkpoint}[An empty overlap is zero]
Why would directly integrating from the larger lower bound to
the smaller upper bound give an unphysical result?
\RuledLines{2}
\end{checkpoint}

\HandoutSection{7}{Hydrogen: scatter out of group 2}
Use the same three-group bounds, ${}^1$H ($\alpha=0$),
$\sigma_s(E)=\sigma_0=20\unit{b}$, and $\phi(E')\propto1/E'$
within incident group 2. Write
\[
 L=\ln(E_1/E_2)=\ln(10^5).
\]
\begin{boardbox}{Set up the two nonzero integrals}
\begin{align*}
 \sigma_{s,2\gets2}
 &=\frac{\sigma_0}{L}\int_{E_2}^{E_1}
                   \frac{E'-E_2}{(E')^2}\dd E',\\
 \sigma_{s,3\gets2}
 &=\frac{\sigma_0}{L}\int_{E_2}^{E_1}
                   \frac{E_2-E_3}{(E')^2}\dd E'.
\end{align*}
Explain the numerators using the allowed outgoing energy intervals.
\RuledLines{2}
\end{boardbox}

\switchcolumn
\RaggedRight
\HandoutSection{8}{Evaluate the scattering column}
The group-to-group cross sections are
\begin{align*}
 \sigma_{s,1\gets2}&=0,\\
 \sigma_{s,2\gets2}&=\sigma_0\frac{L+E_2/E_1-1}{L},\\
 \sigma_{s,3\gets2}&=\sigma_0\frac{(E_2-E_3)(1/E_2-1/E_1)}{L}.
\end{align*}
\begin{boardbox}{Compute the values}
\renewcommand{\arraystretch}{1.6}
\begin{tabularx}{\linewidth}{@{}l>{\centering\arraybackslash}X@{}}
\toprule
Destination & $\sigma_{s,g\gets2}$ (b)\\
\midrule
$g=1$ (fast) & $0$\\
$g=2$ (intermediate) & \MathBlank[1.10in]{}\\
$g=3$ (thermal) & \MathBlank[1.10in]{}\\
\bottomrule
\end{tabularx}
\end{boardbox}
Most collisions stay within the very broad intermediate group;
about $8.68\%$ enter the retained thermal interval.
These are conditional collision fractions, \emph{not} ratios of
steady-state group fluxes.

\begin{checkpoint}[Where did the missing cross section go?]
With $E_3=10^{-3}$ eV, some hydrogen collisions produce
\mbox{$0<E<E_3$}. Their contribution is
\[
 \sigma_{s,<E_3\gets2}
 =\sigma_0\frac{E_3(1/E_2-1/E_1)}{L}
 =\MathBlank[0.95in]{}\unit{b}.
\]
Hence
\[
 \sum_{g=1}^3\sigma_{s,g\gets2}
 =\sigma_0-\sigma_{s,<E_3\gets2}\approx19.99826\unit{b}.
\]
How would the column sum change if the last group extended to $E=0$?
\RuledLines{2}
\end{checkpoint}

\HandoutSection{9}{Check the model, not just the algebra}
For a complete set of outgoing energies,
\[
 \sum_g\Sigma_{s,g\gets g'}=\Sigma_{s,g'},
 \qquad R_x=\sum_g\Sigma_{x,g}\phi_g.
\]
For truncated groups, keep any out-of-range transfer as an
additional loss, or state the approximation used to close the range.
Do not silently claim exact conservation after discarding that transfer.
\begin{takeawaybox}
Multigroup balance determines the \emph{amount} of flux in each
group. The within-group spectra used to prepare cross sections
remain modeling inputs. Fission couples the groups again;
Lesson 14 uses that coupling to construct the $k$-eigenvalue problem.
\end{takeawaybox}
\end{handoutcolumns}
\end{document}
